\begin{frame}<1-4>[t]{例题讲解}
\begin{campuscanvas}
% Source slide 20, objects 2; visible 2-4
\only<2-4|handout:1>{\node[anchor=north west,inner sep=0,text=black] at (70.000,155.000) {\begin{minipage}[t]{98.750mm}\raggedright\fontsize{12}{18.00}\selectfont \textbf{例 5}\quad 用区间表示下列集合：\par （1）$\{x\in\mathbf R\mid -2\le x\le4\}$；\par （2）$\{t\in\mathbf R\mid t<a,\ a\in\mathbf R\}$；\par （3）$\{u\in\mathbf R\mid u\ge0\}$。\end{minipage}};}
% Source slide 20, objects 4; visible 3-4
\only<3-4|handout:1>{\node[anchor=north west,inner sep=0,text=black] at (70.000,420.000) {\begin{minipage}[t]{142.500mm}\raggedright\fontsize{12}{18.00}\selectfont \red{解}\quad （1）$\{x\in\mathbf R\mid -2\le x\le4\}=[-2,4]$；\par （2）$\{t\in\mathbf R\mid t<a,\ a\in\mathbf R\}=(-\infty,a)$；\par （3）$\{u\in\mathbf R\mid u\ge0\}=[0,+\infty)$。\end{minipage}};}
% Source slide 20, objects 32; visible 4
\only<4|handout:1>{\node[anchor=north west,inner sep=0,text=black] at (915.000,155.000) {\begin{minipage}[t]{37.500mm}\raggedright\fontsize{9.5}{14.25}\selectfont \textcolor{campusgreen}{\textbf{约定}}\par 如果从上下文的关系看，$x\in\mathbf R$ 是明确的，那么集合 $\{x\in\mathbf R\mid x<4\}$ 可简写为 $\{x\mid x<4\}$。\end{minipage}};}
\end{campuscanvas}
\end{frame}
